% Makale içi ek (MVA, tek PDF): hakemin görmesi gereken kayıt ve hesaplar. Geri kalan: genişletilmiş ek (ek.tex).
\section{Deney kaydı özeti ve hesap notları}\label{sec:deney-kaydi}

Bu ek, hakemin denetlemesi gereken kayıtları ve hesapları verir: koşum dökümünün özeti, durma kuralının uygulanışı, kilit sonrası işler, GPU bütçesi, ön-eleme olasılıklarının hesabı, gradyan normu oranı ve makullük çıpasının kaynakları. Koşum dökümü tablosu, ek sonuçlar ve yöntem ayrıntıları genişletilmiş ektedir (GitHub deposu, \texttt{paper/ek.pdf}).

Genişletilmiş ekteki Tablo~\ref{tab:deney-kaydi} taban model makullük denetimi, üç tur, kilitli test aşaması ve test sonrası ek tohum koşumlarını kart başına döker. Toplam 60 koşum damgasının beşi çıkış kodu 0 ile bitmedi: üçü kalibrasyon (rc=98), biri yetim süreç (rc=96), biri H05 arızasıdır (rc=95). Kanıt niteliğinde 35 koşum vardır: 12 aşama-3 eğitim, 1 H07 değerlendirmesi ve 9 kilitli test değerlendirmesi (22), ek tohum için 12 eğitim ve 1 WBF tanı koşumu. Kalan 25 koşum kapı, kontrol, kalibrasyon ve arıza kaydıdır ve sonuç olarak raporlanmadı. Toplam 75.45 GPU saati harcandı (hedef 75, tavan 90). Her koşumun tam komutu ve kural parametreleri koşum dizinindeki \texttt{config.json} dosyasında, kuyruklar \texttt{kuyruk/} altında saklıdır.

\subsection{Durma kuralının uygulanışı}\label{sec:ek-durma}
Genişletilmiş ekteki Şekil~\ref{fig:f4-butce}a tur başına GPU süresini ve kümülatif toplamı, (b) tur başına en yüksek eğitim-zamanı kartının $\Delta$'sını gösterir; dolu daire üç tohum ortalamasıdır, boş daire tohum-0 ön-eleme değeridir. Durma kuralı (ardışık iki turda en yüksek eğitim kartı $\Delta<+0.30$ ise kilitli teste geç) ikinci turun sonunda, o tarihte eldeki değerlerle tetiklendi. Birinci turun en yüksek eğitim kartı H03'tü ve üç tohum ortalaması (+0.25) o sırada biliniyordu; ikinci turda yalnız tohum-0 ön-eleme koşumları vardı ve en yüksek kart H04'tü ($-$0.30). Kümülatif süre o noktada 35.30 saatti. Şekildeki ikinci ve üçüncü tur ortalamaları kilit sonrası ek tohumlarla hesaplanmıştır ve durma kararına girmedi.

Kural metni $\Delta$'nın tohum-0 değeri mi yoksa üç tohum ortalaması mı olduğunu belirtmez; her tur için o anda eldeki en iyi tahmin kullanıldı: birinci turda üç tohum ortalaması (+0.25), ikinci turda tohum-0 değeri ($-$0.30). İki tur da tohum-0 değeriyle okunsaydı birinci turun değeri H03'ün tohum-0 farkı (+0.37, yani $\ge+0.30$) olurdu ve kural bu noktada tetiklenmezdi. Kural metni eğitim-zamanı ile eğitimsiz kartı da ayırmaz; uygulamada ikinci turdaki eğitimsiz TTA satırı (doğrulamada, tohum 0, +0.80, $\times$1.10 göreli gecikme kısıtı dışında) kural dışı tutuldu. TTA da tur-2 kartı sayılsaydı kural tetiklenmezdi (+0.80 $\ge$ +0.30); durma kararı bu yüzden eğitim-zamanı yorumuna dayanır.

Doz sorusu için planlanmış tek kart H08, kural tetiklendikten sonra üçüncü turda koşuldu (ön-eleme +0.22, 3.16 saat; üç tohum ortalaması $-$0.08), ardından kilitli teste geçildi (0.21 saat). H08'in tasarımı, kodu ve koşum beklentisi tetiklenmeden yaklaşık 9.6 saat önce, ikinci turun hiçbir sonucu gelmeden yazılmıştı. Kartın kendi dosya damgası bunun kanıtı olarak kullanılamaz, çünkü kart dosyası 2026-10-03'te yeniden yazılmıştır; sıra kuyruk ve kod dosyalarının damgalarından kuruldu. Tur sınırı beş olan kural belgesinde üçüncü tur için bir aday havuzu vardı; ancak kural tetiklendiğinde yine de bir tur koşulmasına izin veren önceden yazılmış bir istisna yoktu. H08'in koşulması bu yüzden önceden yazılmış kurallardan bir sapmadır (Bölüm~\ref{sec:sinirliliklar}); kart kabul almadı ve sapma hiçbir kararı olumluya çevirmedi.

\subsection{Kilit sonrası işler}\label{sec:ek-kilit}
VOC'ta kilit açıldıktan sonra yapılanların tam listesi şudur: (i) ön-elemede elenen altı kartın (H01, H02, H04, H05, H06, H08) tohum 1 ve 2 doğrulama koşumları, yani 12 eğitim koşumu; (ii) H07'nin WBF dalı için doğrulama kümesinde, WBF'nin eşik, atlama eşiği ve füzyon türü çevresinde keşifsel bir tanı ızgarası (hiçbiri teste uygulanmadı; ızgaranın en iyisi kanıt sayılmaz); (iii) n=3, güç ve şans analizleri (mevcut çıktılar üzerinde, yeni koşum yok); (iv) ikinci bir kural belgesiyle (Bölüm~\ref{sec:h05-mekanizma}) H05'in iki bileşen kartı (H05-pos, H05-neg) üç tohumla yalnız doğrulamada koşuldu (6 eğitim koşumu; VOC testi açılmadı); (v) üçüncü bir kural belgesiyle (2026-10-22) H05'in $\lambda_{\mathrm{cls}}$ kazancı telafisi kartı H05c, üç tohumla yalnız doğrulamada koşuldu (3 eğitim koşumu; VOC testi açılmadı). Liste VOC verisindeki işleri kapsar; KITTI deneyi ayrı bir veri kümesinde ve ayrı bir kural belgesiyle yürütülmüştür (Bölüm~\ref{sec:genelleme}). KITTI testi açıldıktan sonra yeni eğitim, test çıkarımı ya da hiperparametre araması yapılmadı; önceden yazılmış parçalar (48 koşumun test değerlendirmesi, H07 test satırları, birincil analiz) dışındakiler keşifsel ya da ikincildir: doğrulamada önceden yazılmış ikincil analiz ve \texttt{best.pt} ile \texttt{last.pt}'nin yeniden değerlendirmesi (96 değerlendirme); mevcut test tahminlerinde kümeli bootstrap $p$, Holm-10, Holm-14 ve işaret-çevirme duyarlılıkları, \texttt{DontCare}'siz $\Delta$, sınıf alt kümeleri ve ikincil tablolar; hiyerarşik bootstrap, sahne grubu başına $\Delta$ ve tram duyarlılığı; dış kaynak değeri; taban modelin 1000 eğitim görüntüsünde tanısı (genişletilmiş ek, Bölüm~\ref{sec:ek-ikincil}).Bunlar dışında VOC'ta başka koşum, hipotez ya da hiperparametre araması yapılmadı ve teste dokunulmadı. Ek tohumlar, n<3 sayıların raporlanmaması ilkesi uyarınca kilit sonrası eklendi ve kilit öncesi eleme kararlarına göre yalnız elenen altı karta verildi.

\subsection{GPU bütçesi}\label{sec:ek-butce}
Kilit açıldıktan sonra altı kartın ek tohum koşumları ve WBF tanı ızgarası 36.79 saat tuttu ve toplamı 38.67 saatten 75.45 saate çıkardı (toplamlar yuvarlanmamış değerlerden hesaplanmıştır); bu, hedefin (75) hafif üstünde, tavanın (90) altındadır.

Genelleme deneyi (Bölüm~\ref{sec:genelleme}) Tablo~\ref{tab:deney-kaydi}'e dahil değildir ve ayrı kaydedilmiştir: 54 tam koşum (KITTI'de iki model $\times$ sekiz kol $\times$ üç tohum $=48$; VOC'ta H05-pos ve H05-neg $\times$ üç tohum $=6$), hepsi 100 epoch. Toplam eğitim süresi 89.02 saattir (YOLO11s 41.15, YOLO26n 28.99, VOC mekanizma kartları 18.88); H05c'nin üç koşumu (VOC, doğrulama; 10.25 GPU saati) ayrıca kaydedilmiştir ve bu toplama dahil değildir. Duman testi, kısa ön deneme, H07 ve test değerlendirmesi de dahil değildir; bu nedenle toplam, ilk deneyin 75.45 saatlik toplamıyla (tüm koşum damgaları) aynı tanımda değildir. KITTI test bölümü 2026-10-22 saat 12:02:15'te, 48 koşumun \texttt{best.pt} ve \texttt{last.pt} ağırlıklarıyla tek açılışta açıldı; VOC testi bu deneyde açılmadı.

\subsection{Ön-eleme olasılıklarının hesabı}\label{sec:ek-onele}
Bölüm~\ref{sec:bulgu-ablasyon}'deki iki olasılık farklı soruları yanıtlar ve yalnız betimseldir. (A) Bilgisiz bir tohum-0 kapısının ($\Delta_{k,0}\ge+0.30$) gerçek etkisi sıfır olan altı kartı da eleme olasılığı $(1-p)^6$ ile hesaplandı; $p=1-\Phi(0.30/\sigma)$ tek kartın kapıdan geçme olasılığıdır ve tek bir ortak $\sigma$ kullanılır. $\sigma$ için iki değer alındı: H03'ün test bootstrap standart hatasından doğrulama boyutuna ölçeklenen görüntü-örnekleme tahmini (0.502; $p=0.275$) ve H03'ün doğrulama tohum-farkı s.d.'si (0.240; $p=0.105$); bunlar sırasıyla 0.145 ve 0.512 verir. (B) ``Elenen kart üç tohumda da kabul almaz'' tanımıyla 0.774, Monte Carlo ile bulundu: kart başına 400\,000 çekim, her çekimde üç tohum $\Delta_s\sim N(0,s_{\Delta,k}^2)$ ($s_{\Delta,k}$ kartın gözlenen tohum-farkı s.d.'si), gerçek etki tam sıfır, kabul ortalama $\Delta\ge+0.40$ ve tek yönlü $t>1.886$. Kartlar bağımsız çekildi; ortak taban modelden doğan korelasyon ve $n=3$ ortalamasının eleme tohumunu içermesi modellenmedi. Hiçbir hesap kapının bilgilendirici olduğunu ya da olmadığını kanıtlamaz; kapının duyarlılığı gerçek pozitif görülmediğinden ölçülemez.

\subsection{Gradyan normu oranı}\label{sec:ek-gradyan}
H05 için raporlanan VFL/BCE oranı, sınıflandırma logitlerine göre alınan kayıp gradyanının L2 normundan hesaplanır: aynı logit ve hedefler için VFL ve BCE gradyan normları her 200. eğitim adımında ölçülür (koşum başına 486 ölçüm) ve norm toplamlarının oranı alınır. H05c'nin manipülasyon denetimindeki etkin gradyan oranı, bu oranın $1.47/0.5$ ile çarpımıdır (Bölüm~\ref{sec:h05}).

\subsection{Makullük çıpasının kaynakları}\label{sec:ek-cipa}
Bölüm~\ref{sec:replikasyon}'deki çıpanın alt ucunu veren iki YOLO11n kaynağından ADL-YOLO sıfırdan eğitimi bildirir, YOLO-Master başlatmayı belirtmez \cite{zhang_2025_adlyolo,lin_2025_yolomaster}. Üst uçtaki YOLOv8n değeri (63.5) kaynağın 4.4. bölümünün metninde doğrulandı, ilgili tablo hücresi okunamadı; kaynak VOC07 testini doğrulama ya da test kümesi olarak kullandığından kontrol noktası muhtemelen testte seçilmiştir \cite{elhamdouni_2025_yolov8nvoc}. Sıfırdan eğitilen YOLOv8n'in 60.2'lik değeri çıpa dışındadır ve yalnız bağlam olarak listelenir (Tablo~\ref{tab:replikasyon}).
